\section{Gauging and anomalies}
\label{sec:gauging}

Throughout this section, $G$ is an affine algebraic group of finite type over a field of
characteristic zero, and $\Rep\left(G\right) = \QC\left(BG\right)$.

Sections \ref{sec:shvcat}--\ref{sec:gauging_anomaly} work in complete generality: $A$ is an
arbitrary $\EE_2$-algebra in $\Pre$, and $G$ acts on it by $\EE_2$-functors. We explain how such
an action produces an object $A^G$ of $\umor_2\left(\Mod_{\Rep\left(G\right)}\left(\Pre\right)\right)$,
and what this gives when $A$ is invertible. In Section \ref{sec:gauging_lines} we specialize to the
case where $A$ is the category of line operators of a $3$-dimensional TQFT, and show that actions
induced by an action on the underlying stack are anomaly-free. In Section \ref{sec:tori} we show that the categories of representations
of quantum groups of tori at roots of unity are invertible relative to quantum Frobenius.
Finally, in Section \ref{sec:aggregated} we discuss an analytic setting.
We begin by recalling sheaves of categories and $1$-affineness in Section \ref{sec:shvcat}.

\subsection{Sheaves of categories and \texorpdfstring{$1$}{1}-affineness}
\label{sec:shvcat}

An ordinary quasi-coherent sheaf on $X$ can be thought of as a functorial assignment
(for every affine $\Spec \left(A\right)$ over $X$) of an $A$-module.
Similarly, a sheaf of categories is a functorial assignment, for every affine $S \to X$,
of a certain module category (in $\Pre$) over the symmetric-monoidal category
$\QC\left(S\right)$.

Formally, $\ShvCat\left(X\right)$ is the limit of $\Mod_{\QC\left(S\right)}\left(\Pre\right)$ over affines
$S\to X$, and $X$ is \emph{$1$-affine} if taking global sections
$\ShvCat\left(X\right)\to\Mod_{\QC\left(X\right)}\left(\Pre\right)$ is an equivalence \cite{G:ShvCat}.
For example, $BG$ is $1$-affine \cite[Theorem 2.2.2]{G:ShvCat}:
$\ShvCat\left(BG\right)\simeq\Mod_{\Rep\left(G\right)}\left(\Pre\right)$.

\subsection{Equivariantization}
\label{sec:gauging_general}

Let $A\in\umor_2\left(\Pre\right)$ be arbitrary.
A sheaf of categories on $BG$ is a $\QC\left(S\right)$-linear category for every
$S\to BG$, compatible with base change.
Computing on the \v{C}ech nerve of $\pt\to BG$, this amounts to a category with a coaction of
$\QC\left(G\right)$.

\begin{defn}
An \emph{action of $G$ on $A$ by $\EE_2$-functors} is an $\EE_2$-algebra $A/G$ in
$\ShvCat\left(BG\right)$ with fiber $A$ at $\pt\to BG$.
Equivalently, it is an $\EE_2$-monoidal coaction
$\rho\colon A\to\QC\left(G\right)\otimes A$, coassociative up to coherent homotopy.
\end{defn}

We write $A/G$ for the sheaf of categories, by analogy with the quotient stack: For a
$G$-stack $X$, the action of $G$ on $\QC\left(X\right)$ is the sheaf of categories on $BG$
whose global sections are $\QC\left(X/G\right)$.
An action is a family of autoequivalences parametrized by the scheme $G$, not merely a
homomorphism from the discrete group $G\left(k\right)$.

\begin{defn}
The \emph{equivariantization} of the action is $A^G\ceqq\Gamma\left(BG,A/G\right)$.
This is the category of objects $a\in A$ equipped with an equivalence
$\rho\left(a\right)\simeq\cO_G\boxtimes a$ satisfying the cocycle condition.
\end{defn}

The functor $\Gamma\left(BG,-\right)\colon\ShvCat\left(BG\right)\to\Mod_{\Rep\left(G\right)}\left(\Pre\right)$ is right
adjoint to the symmetric monoidal functor
$M\mapsto\left(\QC\left(S\right)\otimes_{\Rep\left(G\right)}M\right)_{S\to BG}$.
It is therefore lax symmetric monoidal, and so it sends $\EE_2$-algebras to $\EE_2$-algebras.
Hence $A^G\in\umor_2\left(\Mod_{\Rep\left(G\right)}\left(\Pre\right)\right)$.
Concretely, equivariant objects are multiplied and braided as in $A$.
This is possible because $\rho$ is braided monoidal.
The unit $\Rep\left(G\right)\to A^G$ sends $V$ to $\one_A$ with equivariant structure $V$.
These objects braid trivially with everything, recovering
$\Rep\left(G\right)\subseteq \mathcal{Z}_2\left(A^G\right)$ from the finite setting \cite{DGNO}.

\begin{prop}
\label{prop:equiv}
\leavevmode
\begin{enumerate}
\item Equivariantization defines an equivalence
\begin{equation*}
\umor_2\left(\ShvCat\left(BG\right)\right) \lto{\sim}
\umor_2\left(\Mod_{\Rep\left(G\right)}\left(\Pre\right)\right)\ .
\end{equation*}
Its inverse sends $A'$ to
$\QC\left(\pt\right)\otimes_{\Rep\left(G\right)}A'$ with its induced coaction.
\item $A$ is invertible in $\umor_2\left(\Pre\right)$ if and only if $A^G$ is invertible in
\begin{equation*}
\umor_2\left(\Mod_{\Rep\left(G\right)}\left(\Pre\right)\right).
\end{equation*}
\end{enumerate}
\end{prop}

\begin{proof}
For (1): By $1$-affineness of $BG$ \cite[Theorem 2.2.2]{G:ShvCat},
$\Gamma\left(BG,-\right)$ is a symmetric monoidal equivalence.
The inverse is identified using
\begin{equation*}
\QC\left(\pt\right)\otimes_{\Rep\left(G\right)}\QC\left(\pt\right)\simeq\QC\left(G\right)
\end{equation*}
(Theorem \ref{thm:BZFN}).
This part holds for every $\EE_2$-algebra $A$.

For (2), first assume $A$ is invertible.
By (1), it suffices to show that $A/G$ is invertible in
$\umor_2\left(\ShvCat\left(BG\right)\right)$.
We check the criterion of Theorem \ref{thm:inv_E2} in
$\coef = \ShvCat\left(BG\right)$, using descent: By \cite[Theorem 1.5.7]{G:ShvCat},
\begin{equation*}
\ShvCat\left(BG\right)\simeq\mathrm{Tot}\left(\Mod_{\QC\left(G^\bullet\right)}\left(\Pre\right)\right)
\end{equation*}
as symmetric monoidal $\infty$-categories, with transition functors given by base change,
which are symmetric monoidal and preserve colimits. Write $p_n^*$ for the projections. They
commute with tensor products, relative tensor products and factorization homology, so the
images under $p_n^*$ of $A/G$, $\left(A/G\right)^e$, $\HC\left(A/G\right)$ and their modules
are the base changes $\QC\left(G^n\right)\otimes\left(-\right)$ of the corresponding objects for
$A$.
\begin{itemize}
\item \emph{$3$-dualizability.} An object of a limit of symmetric monoidal $\infty$-categories
is dualizable if and only if all its images are \cite[Lemma 1.4.6]{G:ShvCat}, since the space of
duality data is the limit of the spaces of duality data. The same argument applies to
dualizability of modules. Base change preserves duals, so the dualizability conditions for
$A$ imply those for $A/G$.
\item \emph{The three maps.} Once the relevant modules are dualizable, the centers
$\mathcal{Z}_j\left(A/G\right)$ are computed by duals and relative tensor products, hence commute with
$p_n^*$. A map in a limit of $\infty$-categories is an equivalence if and only if all its images
are, and the images are base changes of the corresponding equivalences for $A$.
\end{itemize}
Conversely, if $A^G$ is invertible, then so is $A/G$ by (1). Taking the fiber at
$\pt\to BG$ is a functor from $\ShvCat\left(BG\right)$ to $\Pre$ which is symmetric monoidal and preserves colimits, so it
induces a symmetric monoidal functor on Morita categories, and sends the invertible $A/G$ to $A$.
\end{proof}

\begin{rmk}
\label{rmk:inv_descent}
The proof of (2) shows in general that an $\EE_2$-algebra in $\ShvCat\left(BG\right)$ is
invertible if and only if its fiber at $\pt\to BG$ is invertible: All conditions of Theorem
\ref{thm:inv_E2} can be checked after pullback along $\pt\to BG$.
\end{rmk}

\subsection{The anomaly of an action on an invertible algebra}
\label{sec:gauging_anomaly}

We remain in the general setting, but now assume that $A$ is invertible, so that $A^G$ is
invertible by Proposition \ref{prop:equiv}(2). We write $\alpha_A$ and $\alpha_{A^G}$ for the
corresponding invertible $4$-dimensional theories \cite{BJSS}.

\begin{nota}
\label{nota:Pic}
For a symmetric monoidal higher category $\coef$, we write $\mathrm{Pic}\left(\coef\right)$ for its
Picard $\infty$-groupoid: the invertible objects of $\coef$, and the invertible morphisms between
them.
\end{nota}

\begin{defn}
\label{defn:Witt}
Let $\coef$ be a symmetric monoidal $\infty$-category. The \emph{Witt group} of $\coef$ is the
group
\begin{equation*}
\mathrm{W}\left(\coef\right)\ceqq\pi_0\mathrm{Pic}\left(\umor_2\left(\coef\right)\right)
\end{equation*}
of Morita classes of invertible $\EE_2$-algebras in $\coef$, under $\otimes$. We use this for
$\coef$ equal to $\Pre$, $\ShvCat\left(BG\right)$ and $\Mod_{\Rep\left(G\right)}\left(\Pre\right)$.
\end{defn}

The inverse of $A$ is $A^{\sigma\rev}$, so we may form the invertible
$\Rep\left(G\right)$-linear $\EE_2$-algebra
\begin{equation*}
\omega\left(A/G\right)\ceqq
A^G\otimes_{\Rep\left(G\right)}\left(\Rep\left(G\right)\otimes A^{\sigma\rev}\right) \ ,
\end{equation*}
which compares $A^G$ with $A$ carrying the trivial action.
Its fiber $A\otimes A^{\sigma\rev}\simeq\one$ is Morita trivial.

\begin{defn}
\label{defn:anomaly}
The \emph{anomaly} of the action is the class
\begin{equation*}
\left[\omega\left(A/G\right)\right]
=\left[A^G\right]-\left[\Rep\left(G\right)\otimes A\right]
\in\mathrm{W}\left(\Mod_{\Rep\left(G\right)}\left(\Pre\right)\right)
\simeq\mathrm{W}\left(\ShvCat\left(BG\right)\right) \ ,
\end{equation*}
where the equivalence is Proposition \ref{prop:equiv}(1).
\end{defn}

The anomaly vanishes exactly when $A^G$ is Morita equivalent over $\Rep\left(G\right)$ to
$\Rep\left(G\right)\otimes A$.

\subsection{Line operators of a \texorpdfstring{$3$}{3}-dimensional theory}
\label{sec:gauging_lines}

We now specialize to line operators (Section \ref{sec:line_ops}).
Let $\mathcal{F}$ be a $3$-dimensional TQFT valued in $\umor_1\left(\Pre\right)$, with $B\ceqq \mathcal{F}\left(\pt\right)$
and category of line operators $A\ceqq \mathcal{F}\left(S^1\right)\simeq \mathcal{Z}_1\left(B\right)$.

A \emph{$G$-symmetry} of $\mathcal{F}$ is a lift of $\mathcal{F}$ to a TQFT $\mathcal{F}_{BG}$ valued in
$\umor_1\left(\ShvCat\left(BG\right)\right)$, i.e.\ a family of theories over $BG$.
On the circle, $\mathcal{F}_{BG}\left(S^1\right)$ is an action of $G$ on $A$ by $\EE_2$-functors,
so the constructions of Sections \ref{sec:gauging_general} and \ref{sec:gauging_anomaly}
apply.
Composing with the equivalence $\Gamma\left(BG,-\right)$ gives a $\Rep\left(G\right)$-linear
TQFT $\mathcal{F}^G$ with
\begin{equation*}
\mathcal{F}^G\left(\pt\right) = B^G \ ,
\qquad
\mathcal{F}^G\left(S^1\right) = A^G \ ,
\end{equation*}
and $B^G$ is a $1$-morphism $A^G\to\one$ in $\umor_2\left(\Mod_{\Rep\left(G\right)}\left(\Pre\right)\right)$.
When $A$ is invertible, $\mathcal{F}^G$ is a boundary condition for $\alpha_{A^G}$.
For example, an action of $G$ on $X$ gives $\QC\left(X\right)^G = \QC\left(X/G\right)$
and $\QC\left(\cL X\right)^G = \QC\left(\cL X/G\right)$.
Here $\cL X/G$ is the relative loop stack of $X/G\to BG$, with relative convolution.

\begin{lem}[Geometric actions are anomaly-free]
\label{lem:geometric}
Let $G$ act on a perfect stack $X$ satisfying diagonal descent (Definition
\ref{defn:diag_descent}), with $X/G$ perfect, so
that $A=\left(\QC\left(\cL X\right),*\right)$ is Morita trivial (Theorem \ref{thm:trivial_QC}). Then
$\QC\left(X/G\right)$ is an invertible $1$-morphism from
$\QC\left(\cL X/G\right)\simeq A^G$ to $\Rep\left(G\right)$ in $\umor_2\left(\Mod_{\Rep\left(G\right)}\left(\Pre\right)\right)$,
and the anomaly of the action of $G$ on $A$ vanishes.
\end{lem}

\begin{proof}
Write $\cL X/G\simeq\left(X/G\right)\times_{\left(X\times X\right)/G}\left(X/G\right)$ for the
relative loop stack of $X/G\to BG$. Theorem \ref{thm:BZFN} for the diagonal of $X/G$ over $BG$
identifies the Drinfeld center of the commutative algebra $\QC\left(X/G\right)$ in
$\Mod_{\Rep\left(G\right)}\left(\Pre\right)$ with $\QC\left(\cL X/G\right)\simeq A^G$. So
$\QC\left(X/G\right)$ is a $1$-morphism $A^G\to\Rep\left(G\right)$ in
$\umor_2\left(\Mod_{\Rep\left(G\right)}\left(\Pre\right)\right)$. Its fiber at $\pt\to BG$ is
$\QC\left(X\right)\colon\QC\left(\cL X\right)\to\one$, since
$\QC\left(\pt\right)\otimes_{\Rep\left(G\right)}\QC\left(Y/G\right)\simeq\QC\left(Y\right)$ for
$Y=X,\cL X$ (Theorem \ref{thm:BZFN}). This fiber is invertible by Theorem \ref{thm:trivial_QC},
and invertibility can be checked after pullback along $\pt\to BG$: As in the proof of
Proposition \ref{prop:equiv}(2), the dualizability conditions and the four maps of Theorem
\ref{thm:trivial_Alg2} are computed termwise in
$\ShvCat\left(BG\right)\simeq\mathrm{Tot}\left(\Mod_{\QC\left(G^\bullet\right)}\left(\Pre\right)\right)$,
and a map in a limit is an equivalence if all its images are. So $\left[A^G\right]=\left[\Rep\left(G\right)\right]=\left[\Rep\left(G\right)\otimes A\right]$,
and $\left[\omega\left(A/G\right)\right]=0$.
\end{proof}

In terms of the $2$-shifted cotangent bundle $T^*\left[2\right]X$ (cf.\ Remark
\ref{rmk:DMod_center}), Lemma \ref{lem:geometric} says that symmetries of $T^*\left[2\right]X$
induced by symmetries of $X$ are anomaly-free.

\subsection{Quantum groups of tori at roots of unity}
\label{sec:detect}

We now apply the constructions of this section to quantum groups of tori at roots of unity, where
the quantum Frobenius gives a central $\EE_2$-functor from $\Rep\left(H\right)$, for a torus $H$, to
the category of representations, landing in transparent objects. We show that the category of
representations is invertible over $\Rep\left(H\right)$ (Theorem \ref{thm:T_invertible}), so that
the anomaly of the action is defined (Definition \ref{defn:anomaly}).

Throughout this subsection, $k=\CC$, and the group acting by $\EE_2$-functors is denoted $H$.

\label{sec:tori}

Let $T$ be a torus with character lattice $\Lambda$, and let $Q\colon\Lambda\to k^\times$ be a
quadratic form, with double braiding
$b\left(\lambda,\mu\right)=Q\left(\lambda+\mu\right)Q\left(\lambda\right)^{-1}Q\left(\mu\right)^{-1}$
and radical $\Lambda_0$ of $b$. Assume that $Q|_{\Lambda_0}=1$ and that $\Lambda/\Lambda_0$ is
finite, and let $H\ceqq\Hom\left(\Lambda_0,\GG_m\right)$. By \cite{JS:braided}, $Q$ determines a
braided structure on $\Lambda$-graded vector spaces, and we write $\Rep\left(T_q\right)$ for its
derived category, i.e.\ $\Lambda$-graded complexes. The basic example is $T=\GG_m$,
$Q\left(a\right)=q^{a^2}$ with $q$ a primitive $n$-th root of unity, $n$ odd, for which
$\Lambda_0=n\ZZ$.

\begin{theorem}
\label{thm:T_invertible}
$\Rep\left(T_q\right)$ is invertible in $\umor_2\left(\Mod_{\Rep\left(H\right)}\left(\Pre\right)\right)$.
\end{theorem}

\begin{proof}
The inclusion $\Vect_{\Lambda_0}\subseteq\Rep\left(T_q\right)$ is an $\EE_2$-map from
$\Rep\left(H\right)=\Vect_{\Lambda_0}$ with its symmetric monoidal structure, because
$Q|_{\Lambda_0}=1$, and it lands in transparent objects, because $\Lambda_0$ is the radical of $b$. So $\Rep\left(T_q\right)\simeq A^H$ for the fiber
$A\ceqq\Rep\left(T_q\right)\otimes_{\Rep\left(H\right)}\Vect$ (Proposition \ref{prop:equiv}(1)). The
fiber is computed on gradings, and is $\Vect_{\Lambda/\Lambda_0}$ with the induced braiding: a
pointed fusion category, which is non-degenerate because $\Lambda_0$ is the radical of $b$. As a
non-degenerate braided fusion category, it is invertible \cite[Theorem 1.1]{BJSS} (Remark \ref{rmk:abelian}). By Proposition \ref{prop:equiv}(2), $A^H$ is
invertible over $\Rep\left(H\right)$.
\end{proof}
